\section*{Appendix C: Sources and References}

Frameworks and disciplines drawn on throughout, for readers who want to consult the originals directly rather than this report's application of them.

\begin{itemize}
\item Kahneman, D. \& Klein, G. (2009). "Conditions for intuitive expertise: a failure to disagree." \emph{American Psychologist}, 64(6). --- Section 7.
\item Snowden, D. \& Boone, M. (2007). "A Leader's Framework for Decision Making." \emph{Harvard Business Review}. --- the Cynefin framework, Section 7.
\item Sutton, R. (2019). "The Bitter Lesson." --- Section 7.
\item Parasuraman, R., Sheridan, T. B., \& Wickens, C. D. (2000). "A model for types and levels of human interaction with automation." \emph{IEEE Transactions on Systems, Man, and Cybernetics}, 30(3). --- Section 7.
\item Christiano, P. Writing on the verification/generation asymmetry in AI alignment. --- Section 7.
\item Polanyi, M. (1966). \emph{The Tacit Dimension}. --- Section 7.
\item Kim, D., Garg, S., Peng, Z., \& Garg, S. "Correlated Errors in Large Language Models," accepted ICML 2025 (arXiv:2506.07962). --- Section 12.
\item OpenAI and Hugging Face, public incident disclosures, July 2026. --- Section 6.
\item UK AI Security Institute (AISI), cybersecurity evaluation findings, July 2026. --- Section 6.
\item "Terminus-4B: Can a Smaller Model Replace Frontier LLMs at Agentic Execution Tasks?" (arXiv:2605.03195). --- Section 8.
\item NIST, ARIA program (model testing, red-teaming, and field testing for AI applications). --- Section 14.
\item Leveson, N. --- STAMP (systems-theoretic accident model and processes). --- Section 11.
\item Failure modes and effects analysis (FMEA) and FMECA, as established reliability-engineering practice rather than a single citable source. --- Sections 11 and 13.
\end{itemize}

All empirical and benchmark-specific claims in Section 8 should be treated as dated regardless of what's listed above; re-verify against current sources before relying on them.
